\documentclass[11pt,a4paper]{article}
\usepackage[margin=2.1cm]{geometry}
\usepackage[T1]{fontenc}
\usepackage{times}
\usepackage[version=4]{mhchem}
\usepackage{enumitem}
\usepackage[hidelinks]{hyperref}
\setlength{\parskip}{0.3em}
\setlength{\parindent}{0pt}
\pagestyle{empty}

\begin{document}

\begin{flushright}
Yaming Hu\\
Independent Researcher\\
Guiyang, China\\
Email: \href{mailto:64687555@qq.com}{64687555@qq.com}\\
ORCID: 0009-0003-1406-0485\\[0.4em]
August 17, 2026
\end{flushright}

\vspace{0.3em}

\begin{flushleft}
Editor-in-Chief\\
\emph{Metallurgical and Materials Transactions B}\\
Springer / TMS
\end{flushleft}

\vspace{0.4em}

Dear Editor-in-Chief,

\textbf{Re: Submission of an original research article --- ``Thermodynamic and Kinetic Assessment of Titanium Production via In-Situ Electrolysis and Bimetallic Complementary Reduction in MgCl$_2$--CaCl$_2$--NaCl--KCl Molten Salt''}

I am pleased to submit the above-referenced manuscript (34 pages, 11 figures, 10 tables, 27 references) for consideration as a \emph{research article} in MMTB.

The titanium industry remains dominated by the batch-type Kroll process, with its long reduction cycles, high energy intensity, and limited scalability. Building on the molten-salt electrochemistry of the FFC Cambridge and OS processes, this work proposes and quantitatively assesses an alternative concept: \emph{in-situ} electrolytic co-deposition of Mg and Ca in a low-melting quaternary MgCl$_2$--CaCl$_2$--NaCl--KCl melt (580--650~$^{\circ}$C), in which Mg performs fast bulk reduction of \ce{TiO2} while Ca provides deep deoxidation, regenerating both reductants electrochemically within the same cell.

The main contributions of the manuscript are:

\begin{enumerate}[itemsep=2pt,topsep=2pt]
  \item \textbf{A self-consistent thermodynamic framework}: all reduction pathways are strongly exergonic ($\Delta G^{\circ} = -231.0$ to $-307.5$~kJ/mol at 600~$^{\circ}$C); the Mg/Ca deoxidation limits (3908 vs.\ 20.2~ppm O, 193$\times$) provide a quantitative basis for the bimetallic design.
  \item \textbf{Electrochemical design criteria}: a \ce{Ca^{2+}/Na+} co-deposition window of +120~mV (ideal), positive (74.9--147.1~mV) across the activity-coefficient sensitivity range, plus an anode \ce{O2}/\ce{Cl2} selectivity boundary and Butler--Volmer overpotential analysis.
  \item \textbf{Kinetic and process-energy assessments}: shrinking-core kinetics predict $\sim$6.3~min for complete reduction of 100~$\mu$m particles; the conditional net energy consumption is 11{,}191~kWh/ton-Ti (second-law efficiency $\sim$40.7\%), and a Monte Carlo uncertainty analysis ($P(W<14{,}000) = 85.8\%$) exposes the dominance of the unvalidated current-efficiency parameter.
\end{enumerate}

The paper is deliberately framed as a \emph{theoretical feasibility assessment}: two prerequisites (oxide solubility $\geq$0.5~wt\% and sustained current efficiency $>$0.70 with an inert oxygen-evolving anode) are identified as mandatory experimental validation targets. All calculations are fully reproducible; the Python scripts and data tables are publicly available at\\ \url{https://github.com/huyamingc/ti-in-situ-electrolysis}. A preprint version of this work has been deposited on Zenodo (DOI: \href{https://doi.org/10.5281/zenodo.21973470}{10.5281/zenodo.21973470}).

I believe this work falls squarely within the scope of \emph{Metallurgical and Materials Transactions B}---process metallurgy, molten-salt electrochemistry, and process modeling---and will interest readers working on emerging titanium production routes.

I confirm that this manuscript is original, unpublished, and not under consideration by any other journal; I have no conflicts of interest to disclose and received no funding. As the sole author, I have approved the manuscript and agree with its submission.

\vspace{0.5em}

Sincerely,

\vspace{0.5em}

Yaming Hu\\
Independent Researcher, Guiyang, China

\end{document}
